\section{Introduction, Literature Survey, Methodology, Objectives}

\subsection{Introduction}
\label{sec:introduction}

A supernova is the catastrophic, once-per-object brightening of a star
(or stellar remnant) at the end point of its evolution. For a period of
days to months it can rival or exceed the integrated brightness of its
entire host galaxy, which is what makes supernovae visible across cosmological
distances and useful as tracers of the universe's structure and expansion.
Astronomers group observed supernovae by \emph{spectroscopic type}: what
absorption and emission features appear when the transient's light is
dispersed through a spectrograph. This report is concerned with three of
those types, and with a fourth category -- ``none of the above'' -- that a
real alert stream inevitably also produces.

\paragraph{Two unrelated physical channels, one shared naming scheme.}
Type Ia supernovae arise from the thermonuclear disruption of a carbon-oxygen
white dwarf, either through mass accretion from a companion star past a
critical threshold or through the merger of two white dwarfs. The explosion
consumes the star with no hydrogen involved at any stage, which is why a Ia
spectrum shows no hydrogen lines. Type Ib/c and Type II supernovae, by
contrast, are both \emph{core-collapse} events: a massive star ($\gtrsim
8\,M_\odot$) exhausts its nuclear fuel, its iron core can no longer support
itself against gravity, and the resulting collapse (typically to a neutron
star) drives a shock that unbinds the star's outer layers. Type II
supernovae still have their hydrogen envelope at the moment of collapse;
Type Ib and Ic supernovae have already lost it (and, for Ic, their helium
layer too) to strong stellar winds or to a binary companion. It is worth
stating plainly, because the shared Roman-numeral naming scheme actively
obscures it: \emph{Type Ia is not physically related to Type Ib or Ic}. The
``I'' in each name refers only to the historical absence of hydrogen lines
in the discovery-era spectroscopic classification scheme
\citep{filippenko1997}, not to a shared progenitor or explosion mechanism.
Type Ib and Ic are two flavours of the same core-collapse channel as Type
II; Type Ia is a different channel entirely. This project follows
convention (and its own taxonomy, Section~\ref{sec:data}) in merging Ib and
Ic into one target class, since distinguishing a stripped-envelope
supernova that has lost its hydrogen from one that has lost its helium as
well requires spectroscopy this pipeline does not have.

\paragraph{Why Type Ia matters beyond stellar physics.}
Because a Type Ia explosion detonates a white dwarf at (or close to) a
common physical trigger mass, the resulting explosions are photometrically
similar to one another: their peak absolute brightness is
\emph{standardisable} through empirical relations between peak luminosity,
light-curve decline rate, and colour. That standardisability is what turns
a Type Ia supernova into a ``standard candle'' -- an object whose true
brightness can be inferred well enough from its light curve alone that
comparing it to its \emph{observed} brightness yields a distance. Type Ia
supernovae were the tool used to discover the accelerating expansion of the
universe, and they remain a primary distance-ladder rung today. Recovering
Ia identity from photometry alone, quickly, is therefore not merely a
pattern-recognition exercise: it is a prerequisite for using a given object
cosmologically at all.

\paragraph{Spectroscopy is definitive; photometry is what scales.}
A spectrum directly measures the physics -- the elements present, their
ionisation state, the ejecta's expansion velocity -- and is the only fully
reliable way to assign a supernova's type. It is also expensive: it requires
a sufficiently large telescope, a spectrograph, and dedicated observing time
allocated to one object at a time. What a wide-field synoptic survey
produces automatically, for every object, every time it is observed, is a
\emph{light curve}: a sequence of calibrated flux measurements at a handful
of timestamps and wavelength filters (``bands''). The central question this
project asks is how much of the spectroscopic answer -- Ia, Ib/c, or II --
can be recovered from a light curve's shape alone: how fast it rises, how
it declines, whether it plateaus, all without ever taking a spectrum. Doing
this \emph{early}, before the light curve has finished playing out, is the
useful case in practice: a system that can say ``point a spectrograph at
this one'' while the transient is still spectroscopically interesting is
worth far more than one that can only classify a supernova after it has
faded.

\paragraph{The scale problem.}
The Zwicky Transient Facility (ZTF) already produces alert volumes that
strain manual and even semi-automated triage
\citep{bellm2019ztf}. The Vera C.\ Rubin Observatory's Legacy Survey of
Space and Time (LSST), whose main ten-year survey began operations on 30
June 2026, is designed to generate on the order of $10^7$ transient and
variable-source alerts per night at full operational cadence
\citep{ivezic2019lsst} -- a volume no spectroscopic follow-up programme,
however well resourced, can classify object by object. Automated,
photometry-only triage is therefore not an optional convenience for the
LSST era; it is the only mechanism by which a human-in-the-loop follow-up
programme can be pointed at the right handful of objects out of millions.
This report describes one such triage system, built and measured against
both simulated LSST-like data and real ZTF data, and served, briefly, over
a live alert stream from both surveys.
